\documentclass[12pt]{article}
\usepackage{amsmath,amssymb,amsfonts}
\usepackage[margin=1in]{geometry}

\begin{document}

\section*{Problem 5.25}

\textbf{Problem:} Let us define $E(\lambda)$, an operator-valued function, as in Section 5.11. Denote by $A$ and $\phi_\nu$ ($\nu = 1, 2, \ldots, N$) the nondegenerate eigenvalues and corresponding eigenvectors of an Hermitian operator $A$, and call $P_\nu$ the projection operator associated with $A$. Then define:
\[
E(\lambda) = \begin{cases} 0 & \text{for } \lambda < \lambda_1, \\ \sum_{\nu=1}^{k} P_\nu & \text{for } \lambda_k \leq \lambda < \lambda_{k+1}, \quad \nu = 1, 2, \ldots, N-1, \\ I & \text{for } \lambda \geq \lambda_N. \end{cases}
\]

Show that:
\begin{itemize}
\item[(a)] $E(\lambda)E(\mu) = E(\min(\lambda, \mu))$, where $\lambda_{\min} = \min(\lambda_1, \lambda_2)$.
\item[(b)] $E(\lambda)^2 = E(\lambda)$ (i.e., $E(\lambda)$ is a projection).
\item[(c)] If $\lambda \geq \lambda_N$, $E(\lambda) = I$ (the identity).
\end{itemize}

\bigskip
\textbf{Solution:}

\subsection*{Part (a): $E(\lambda)E(\mu) = E(\min(\lambda,\mu))$}

Without loss of generality, assume $\lambda \leq \mu$. Then $\min(\lambda,\mu) = \lambda$.

If $\lambda_k \leq \lambda < \lambda_{k+1}$ and $\lambda_j \leq \mu < \lambda_{j+1}$, then $k \leq j$ (since $\lambda \leq \mu$), and:
\[
E(\lambda) = \sum_{\nu=1}^{k}P_\nu, \qquad E(\mu) = \sum_{\nu=1}^{j}P_\nu.
\]

Since the projection operators are orthogonal ($P_\nu P_\mu = \delta_{\nu\mu}P_\nu$):
\[
E(\lambda)E(\mu) = \left(\sum_{\nu=1}^{k}P_\nu\right)\left(\sum_{\sigma=1}^{j}P_\sigma\right) = \sum_{\nu=1}^{k}\sum_{\sigma=1}^{j}P_\nu P_\sigma = \sum_{\nu=1}^{k}P_\nu = E(\lambda) = E(\min(\lambda,\mu)).
\]

\[
\boxed{E(\lambda)E(\mu) = E(\min(\lambda,\mu)).}
\]

\subsection*{Part (b): $E(\lambda)^2 = E(\lambda)$}

This is a special case of part (a) with $\mu = \lambda$:
\[
E(\lambda)^2 = E(\lambda)E(\lambda) = E(\min(\lambda,\lambda)) = E(\lambda).
\]

Alternatively, directly:
\[
E(\lambda)^2 = \left(\sum_{\nu=1}^{k}P_\nu\right)^2 = \sum_{\nu=1}^{k}\sum_{\sigma=1}^{k}P_\nu P_\sigma = \sum_{\nu=1}^{k}P_\nu^2 = \sum_{\nu=1}^{k}P_\nu = E(\lambda).
\]

\[
\boxed{E(\lambda)^2 = E(\lambda).}
\]

So $E(\lambda)$ is indeed a projection operator for each value of $\lambda$.

\subsection*{Part (c): $E(\lambda) = I$ for $\lambda \geq \lambda_N$}

For $\lambda \geq \lambda_N$, by definition:
\[
E(\lambda) = \sum_{\nu=1}^{N}P_\nu.
\]

Since $\{P_\nu\}$ are the projection operators for the complete set of eigenvectors, the completeness relation (resolution of the identity) gives:
\[
\sum_{\nu=1}^{N}P_\nu = I.
\]

\[
\boxed{E(\lambda) = I \text{ for } \lambda \geq \lambda_N.}
\]

This follows from the spectral theorem: the identity operator is the sum of all projection operators onto the eigenspaces.

\end{document}
